% Section 6 -- Orthographic counterfactuals. Budget: 1.00 page, holds Table 3 and Figure 4.
% Follows docs/DESIGN_DECISIONS.md sections 6 and 8.6 (per-model ATE, three-arm contrast,
% tone-isolation DiD, placement direction, no mediation language).
% STATUS: hypotheses only. No model has been run. Every result is a \placeholder.
\section{Orthographic Counterfactuals}
\label{sec:counterfactuals}

The arms change only the code points, hence the token sequence, so a difference between arms is caused by tokenization and by nothing else about the item. For arm $a$ on model $m$ the estimand is $\tau_m(a) = \mathbb{E}_i[Y_i(a) - Y_i(\mathrm{base})]$, identified by scoring the same items under both arms with all else fixed. C1 and C1$'$ run on \noilai{} and on the 500 Vietnamese test items of XCOPA \citep{ponti-etal-2020-xcopa}, pure NFC in the older convention; C2 runs on the placement-arm set only, since it touches 43 XCOPA items.

\paragraph{Pre-registered hypotheses.}
\hyp{4} (\emph{the rarer convention costs}) leads, because both placement conventions are common, so distribution shift cannot explain a C2 effect: the convention rarer in a reference corpus (expected: the newer; fixed by a count before any run) lowers accuracy on affected items, claimed as a pooled estimate with the number of models whose Holm-adjusted interval excludes zero and equivalence tests for the rest. \hyp{3} (\emph{a crossover}): in pass-through tokenizers $\tau_m(\mathrm{NFD})$ is positive on \task{T1} \var{3}, where NFD makes the tone an explicit symbol, and negative on XCOPA, where it only lengthens the text. The default explanation it must beat is distribution shift: NFD is rare text, and under Gemma~3 it is longer for 91\% of words, carries 59\% more tokens per syllable and has alignment 0.27, so a uniform drop is reported as out-of-distribution sensitivity. The per-family table is primary, and with fewer than five informative families the per-family intervals decide; normalizing tokenizers are the negative control. \hyp{3b}: in Gemma~3 the NFD gain on \var{3} is larger where the tone mark becomes its own token than on matched items where it stays attached, and the same contrast on \var{1} is near zero. Holm runs over the fourteen cells of each model's row.

\paragraph{Effects.}
\placeholder{Table~\ref{tab:arms} reading; verdicts on \hyp{4}, \hyp{3} and \hyp{3b} from the pooled estimates (43 affected XCOPA items descriptive, Appendix~\ref{app:extra}).}

\begin{table*}[t]
  \centering\small
  \begin{tabular}{@{}llrrrrrrr@{}}
    \toprule
    Model & Census & $\Delta$tok & \multicolumn{2}{c}{C1: NFD} & C1$'$: PC & C2: placement & C3a: no tones & XCOPA C1 \\
    & & C1 & \task{T1} all & \task{T1} \var{3} & \task{T1} all & \task{T1} affected & \task{T1} all & \\
    \midrule
    \placeholder{per model} & \placeholder{verdict} & & & & & & & \\
    \bottomrule
  \end{tabular}
  \caption{Within-item effect of each re-encoding arm on accuracy (percentage points, arm minus baseline), paired over identical items, with 95\% clustered intervals; bold: significant after Holm within the model's row. Census: passes through, normalizes (zero by construction) or corrupts (arm refused). \placeholder{TODO: filled by \texttt{noilai.stats.tables.intervention\_table}.}}
  \label{tab:arms}
\end{table*}

\paragraph{Does the effect track token count or the tone symbol?}
Token count is a coarsening of the intervention, not a mediator, so we do not estimate a mediated share. Four pre-registered analyses (Appendix~\ref{app:analysis}) ask whether count can account for the effect: effect modification by the token-count change (associational), a zero-dose contrast, the ordered contrast baseline $<$ PC $<$ NFD, and encoding $\times$ character spacing. \placeholder{Reading of Figure~\ref{fig:deltatok} (Appendix~\ref{app:extra}).}

